\documentclass[a4paper,11pt]{article}
\usepackage[utf8]{inputenc}
\usepackage[T1]{fontenc}
\usepackage[lithuanian]{babel}
\usepackage{enumitem}
\usepackage{geometry}
\geometry{margin=2.5cm}

\title{Sistemos vizija\\\large Individualių ir grupinių užsiėmimų valdymas}
\begin{document}

\maketitle

\section{Sistemos apžvalga}

Sistema skirta individualių ir grupinių užsiėmimų valdymui bei vykdymui baseinuose ir kitose sporto erdvėse. Ji padės vedėjams iš anksto parengti užsiėmimų programas, valdyti dalyvius, pateikti vaizdinę mokomąją medžiagą ir perduoti garso instrukcijas užsiėmimo metu.

Sprendimas bus realizuojamas kaip žiniatinklio kliento-serverio sistema. Pagrindinė aplikacija ir duomenys bus valdomi centralizuotai, o vedėjai sistemą pasieks per žiniatinklio sąsają. Užsiėmimo metu sistema sąveikaus su sporto erdvėje naudojama išorine garso ir vaizdo įranga.

\section{Esminiai architektūriniai sprendimai}

Numatomi šie preliminarūs architektūriniai sprendimai:

\begin{itemize}
    \item Sistema naudos kliento-serverio architektūrą.
    \item Pagrindinė naudotojo sąsaja bus žiniatinklio pagrindu ir pasiekiama standartiniais kompiuteriais bei planšetėmis.
    \item Sistemos duomenys, užsiėmimų programos, dalyvių informacija ir užsiėmimų istorija bus saugomi centralizuotai.
    \item Vaizdo ir kitas užsiėmimų turinys bus valdomas sistemoje ir rodomas prijungtuose ekranuose ar televizoriuose.
    \item Garso instrukcijos ir vedėjo balsas bus perduodami naudojant konkrečiai sporto erdvei tinkamą garso įrangą, įskaitant dalyvių ausines, kai jos reikalingos.
    \item Architektūra turi leisti tą patį programinį sprendimą pritaikyti skirtingoms sporto erdvėms ir skirtingoms techninės įrangos konfigūracijoms.
\end{itemize}

\section{Sistemos skaidymas į dalis}

Projekto planavimo tikslais sistema skaidoma į šias pagrindines dalis:

\subsection{Užsiėmimų programų valdymas}

Atsakingas už užsiėmimų programų kūrimą, redagavimą, saugojimą ir valdymą. Programa sudaroma iš etapų, pratimų, vaizdo turinio, garso instrukcijų ir kitos užsiėmimui reikalingos informacijos.

\subsection{Užsiėmimų valdymas}

Atsakingas už individualių ir grupinių užsiėmimų planavimą bei vykdymą. Modulis valdo parengtos programos vykdymą ir leidžia vedėjui užsiėmimo metu programą pradėti, pristabdyti, tęsti, praleisti etapą arba pereiti prie kito etapo.

\subsection{Dalyvių valdymas}

Atsakingas už dalyvių informacijos saugojimą, dalyvių priskyrimą užsiėmimams ir jų dalyvavimo istorijos valdymą.

\subsection{Vaizdinio turinio pateikimas}

Atsakingas už pratimų, demonstracijų ir kitos vaizdinės informacijos pateikimą televizoriuje, monitoriuje ar kitame tinkamame ekrane. Tai ypač svarbu baseinuose, kur vedėjas ne visada gali tinkamai pademonstruoti vandenyje ar po vandeniu atliekamus veiksmus.

\subsection{Garso komunikacija}

Atsakinga už iš anksto įrašytų garso instrukcijų ir vedėjo balso perdavimą dalyviams. Sprendimas turi palaikyti konkrečiai sporto erdvei tinkamą garso įrangą, įskaitant dalyvių ausines, kai jos reikalingos.

\subsection{Duomenų ir istorijos valdymas}

Atsakingas už užsiėmimų programų, užsiėmimų, dalyvavimo informacijos ir susijusių istorinių duomenų saugojimą.

\section{Preliminarūs techninės įrangos poreikiai}

Sistemai reikalinga standartinė kompiuterinė infrastruktūra bei tiesiogiai sporto erdvėje naudojama techninė įranga.

\begin{itemize}
    \item Serveris arba debesijos aplinka aplikacijai, duomenų bazei ir užsiėmimų turiniui talpinti.
    \item Kompiuteris arba planšetė, kurią vedėjas naudos sistemai pasiekti ir valdyti.
    \item Televizorius, monitorius arba kitas ekranas vaizdiniam užsiėmimų turiniui rodyti.
    \item Ne baseino aplinkoje gali būti naudojamas mobilus televizorius ar ekranas ant ratukų, jei nėra stacionaraus ekrano.
    \item Garso įranga instrukcijoms ir vedėjo balsui perduoti.
    \item Dalyvių ausinės, kai reikalingas individualus garso perdavimas.
    \item Vedėjo mikrofonas arba ausinės su mikrofonu, kai naudojamas tiesioginis balso perdavimas.
    \item Tinklo ryšys tarp aplikacijos ir užsiėmimo metu naudojamos įrangos.
\end{itemize}

Baseino aplinkoje šalia vandens naudojama įranga turi būti parenkama atsižvelgiant į aplinkos sąlygas ir reikalingą atsparumo vandeniui lygį. Konkretūs techninės įrangos modeliai ir komunikacijos technologijos bus parinkti detalaus sistemos projektavimo ir integracinio testavimo metu.

\section{Galimi sistemos išplėtimai}

Architektūra turi sudaryti galimybę ateityje plėsti sistemą neatliekant esminių pagrindinės sistemos funkcionalumo pakeitimų. Galimi išplėtimai:

\begin{itemize}
    \item specializuotos mobiliosios sąsajos vedėjams ir dalyviams;
    \item dalyvių identifikavimas sporto erdvėje;
    \item pakartotinai naudojamų užsiėmimų programų ir vaizdo turinio biblioteka;
    \item individualizuotos programos pagal dalyvio poreikius, pasirengimo lygį ar ankstesnę veiklą;
    \item dalyvių progreso ir užsiėmimų rezultatų analizė;
    \item sistemos naudojimas keliose sporto erdvėse ar skirtinguose sporto centruose.
\end{itemize}

\end{document}
